%% ---------------------------------------------------------------------------
%% Tree Models report template (CS 4372, HW2)
%%
%% How to use:
%%   1. Edit the "PROJECT METADATA" block below (title, dataset, team).
%%   2. Work through the notebook. Every figure it saves to figs/ is picked up
%%      automatically by \figslot below; until then a grey placeholder box shows.
%%   3. Replace each \todo{...} with your own text. Search for "todo" to find them.
%%   4. Do NOT put code or code snippets in this report (assignment requirement).
%% ---------------------------------------------------------------------------
\documentclass[11pt]{article}

\usepackage[T1]{fontenc}
\usepackage[utf8]{inputenc}
\usepackage{lmodern}
\usepackage{microtype}
\usepackage{amsmath}
\usepackage[margin=1in]{geometry}
\usepackage{setspace}
\usepackage{graphicx}
\usepackage{booktabs}
\usepackage{tabularx}
\usepackage{array}
\usepackage{float}
\usepackage[font=small,labelfont=bf]{caption}
\usepackage{xcolor}
\usepackage{hyperref}
\usepackage{fancyhdr}
\usepackage{enumitem}
\usepackage[numbers,sort&compress]{natbib}

\graphicspath{{figs/}}
\definecolor{accent}{HTML}{1F4E79}

% =========================== PROJECT METADATA ===============================
\newcommand{\coursecode}{CS 4372}
\newcommand{\instructor}{A. Nagar}
\newcommand{\projecttitle}{Tree-Based Models on the TODO Dataset}
\newcommand{\projectsubtitle}{Predicting TODO (target) from TODO (predictors)}
\newcommand{\memberA}{Team Member 1}
\newcommand{\netidA}{abc123456}
\newcommand{\memberB}{Team Member 2}      % delete the \memberB/\netidB lines on the
\newcommand{\netidB}{xyz123456}           % cover page if working alone
\newcommand{\headerauthors}{Member 1 and Member 2}
\newcommand{\datasetname}{TODO dataset name}
\newcommand{\dataseturl}{https://archive.ics.uci.edu/dataset/TODO}
% ============================================================================

% Highlighted reminder; remove every use before submitting.
\DeclareRobustCommand{\todo}[1]{{\color{red}\textit{[TODO: #1]}}}

% Figure slot: shows the figure if figs/<name>.pdf exists, else a grey box.
% Usage: \figslot{file-stem}{width as fraction of textwidth}{caption}{label}
\newcommand{\figslot}[4]{%
  \begin{figure}[htbp]
    \centering
    \IfFileExists{figs/#1.pdf}{%
      \includegraphics[width=#2\textwidth]{#1}%
    }{%
      \fbox{\parbox[c][1.6in][c]{#2\textwidth}{\centering\color{gray}
        Placeholder: run the notebook to create \texttt{figs/\detokenize{#1}.pdf}}}%
    }
    \caption{#3}
    \label{#4}
  \end{figure}%
}

\hypersetup{
  colorlinks=true,
  linkcolor=accent,
  citecolor=accent,
  urlcolor=accent,
  pdftitle={Tree-Based Models Report}
}
\setlength{\parindent}{0pt}
\setlength{\parskip}{0.55em}
\setlength{\tabcolsep}{5pt}
\renewcommand{\arraystretch}{1.12}
\setlist[itemize]{leftmargin=1.5em,itemsep=0.2em,topsep=0.3em}

\pagestyle{fancy}
\setlength{\headheight}{14pt}
\fancyhf{}
\fancyhead[L]{\coursecode: Tree-Based Models}
\fancyhead[R]{\headerauthors}
\fancyfoot[C]{\thepage}
\renewcommand{\headrulewidth}{0.4pt}

\begin{document}

% ------------------------------- Cover page ---------------------------------
\pagenumbering{roman}
\begin{titlepage}
  \thispagestyle{empty}
  \centering
  \vspace*{1.15in}

  {\Large\textcolor{accent}{\coursecode}\par}
  \vspace{0.35in}
  {\Huge\bfseries \projecttitle\par}
  \vspace{0.25in}
  {\Large \projectsubtitle\par}

  \vspace{0.55in}
  \rule{0.78\textwidth}{0.6pt}
  \vspace{0.5in}

  {\large
  \begin{tabular}{>{\bfseries}rl}
    Team member: & \memberA \\
    NetID: & \netidA \\
    \addlinespace[0.3em]
    Team member: & \memberB \\
    NetID: & \netidB \\
  \end{tabular}
  \par}

  \vfill
  {\large The University of Texas at Dallas\par}
  \vspace{0.12in}
  {\large Instructor: \instructor\par}
  \vspace{0.4in}
  {\small Dataset: \datasetname\par}
  {\small\url{\dataseturl}\par}
\end{titlepage}

\pagenumbering{arabic}

% -------------------------------- Abstract ----------------------------------
\begin{abstract}
\todo{Write this last, in 150--200 words. Cover: the prediction task and dataset; the
data-quality actions taken; the four tuned models (decision tree, random forest, AdaBoost,
XGBoost); the best model and its headline test metrics; what the diagnostics (confusion
matrix, ROC, precision--recall) revealed; and the main limitation.}
\end{abstract}

% ============================ 1. OBJECTIVE & DATA ===========================
\section{Project Objective and Dataset}

\todo{State the prediction goal and why it matters. Name the dataset, its source, its size
(rows and columns), and cite it \citep{dataset}. Identify the
target variable and whether the task is classification or regression, and say why you chose
it.}

\todo{Describe the predictors in plain language and list any units or encodings documented by
the source. State any assumptions here (the assignment asks for assumptions to be stated
completely).}

\begin{table}[H]
\centering
\caption{Variables in the \datasetname{} dataset.}
\label{tab:variables}
\small
\begin{tabularx}{0.92\textwidth}{lXl}
\toprule
Variable & Meaning & Type \\
\midrule
\todo{name} & \todo{meaning} & \todo{type} \\
\todo{name} & \todo{meaning} & \todo{type} \\
\todo{target} & \todo{meaning; mark as response} & \todo{type} \\
\bottomrule
\end{tabularx}
\end{table}

% ======================= 2. DATA QUALITY & PREPROCESSING ====================
\section{Data Quality and Preprocessing}

\subsection{Loading and consistency checks}

\todo{Explain how the data were loaded from the public URL. Report missing values,
duplicates, impossible or inconsistent values, and what you did about each. State
whether any rule used the target.}

\begin{table}[H]
\centering
\caption{Data audit and final sample sizes.}
\label{tab:audit}
\begin{tabular}{lr}
\toprule
Audit item & Result \\
\midrule
Raw observations & \todo{n} \\
Missing cells & \todo{n} \\
Duplicate rows & \todo{n} \\
Inconsistent rows removed & \todo{n} \\
Final observations & \todo{n} \\
Training observations & \todo{n} \\
Test observations & \todo{n} \\
\bottomrule
\end{tabular}
\end{table}

\subsection{Distributions and categorical conversion}

\todo{Summarize the attribute distributions. Are they normally distributed? If not, give
plausible reasons. Discuss class balance if the task is classification. Explain how
categorical variables were converted and why that encoding was chosen.}

\figslot{01_distributions}{0.96}{Distributions of the predictors.}{fig:distributions}
\figslot{02_target_and_categoricals}{0.90}{Target distribution and its relationship to
categorical predictors.}{fig:target}

% ================== 3. CORRELATION, SCALING, FEATURE SELECTION ==============
\section{Correlation, Scaling, and Feature Selection}

\todo{Report the strongest predictor--target relationships and the most correlated predictor
pairs (give numbers). Explain what multicollinearity means for each of the four tree models.}

\figslot{03_correlation_heatmap}{0.76}{Training-set correlations among predictors and the
target.}{fig:correlation}

\todo{Describe the standardization and normalization steps, what each was fit on, and
whether tree models actually need them. Then explain how the important attributes were
chosen and why you did not use all attributes blindly.}

\figslot{05_feature_selection}{0.72}{\todo{Describe the feature-selection evidence.}}{fig:selection}

% ====================== 4. MODEL CONSTRUCTION & TUNING =======================
\section{Model Construction and Tuning}

\todo{State the train/test split (proportions, random seed) and the cross-validation scheme
used inside the training set. Name the library \citep{pedregosa2011scikit}. Explain the scoring metric used by GridSearchCV and why.}

\begin{table}[H]
\centering
\caption{Hyperparameter grids searched with \texttt{GridSearchCV}.}
\label{tab:grids}
\small
\begin{tabularx}{\textwidth}{lXr}
\toprule
Model & Parameters and values searched & Combinations \\
\midrule
Decision tree & \todo{e.g.\ depth, min samples, criterion} & \todo{n} \\
Random forest & \todo{...} & \todo{n} \\
AdaBoost & \todo{...} & \todo{n} \\
XGBoost & \todo{...} & \todo{n} \\
\bottomrule
\end{tabularx}
\end{table}

\subsection{Decision tree}
\todo{Introduce the method \citep{breiman1984cart}. Give best parameters, cross-validated score, and why you are convinced these are the best
values (e.g.\ the optimum is interior to the grid, not on its edge).}
\figslot{06_dt_tuning}{0.78}{\todo{Decision tree tuning results.}}{fig:dttuning}
\figslot{07_dt_tree}{0.96}{\todo{The best decision tree.}}{fig:dttree}

\subsection{Random forest}
\todo{Same structure as above \citep{breiman2001random}. Add what the ensemble gains over a single tree.}
\figslot{08_rf_tuning}{0.78}{\todo{Random forest tuning results.}}{fig:rftuning}
\figslot{09_rf_tree_and_importance}{0.94}{\todo{One tree from the best forest, with feature
importances.}}{fig:rftree}

\subsection{AdaBoost}
\todo{Same structure as above \citep{freund1997decision}. Note the base-estimator choice and the learning-rate
trade-off.}
\figslot{10_ada_tuning}{0.78}{\todo{AdaBoost tuning results.}}{fig:adatuning}
\figslot{11_ada_stump_and_importance}{0.94}{\todo{A boosted tree with feature
importances.}}{fig:adatree}

\subsection{XGBoost}
\todo{Same structure as above \citep{chen2016xgboost}. Mention the regularization parameters you tuned.}
\figslot{12_xgb_tuning}{0.78}{\todo{XGBoost tuning results.}}{fig:xgbtuning}
\figslot{13_xgb_tree_and_importance}{0.94}{\todo{An XGBoost tree with feature
importances.}}{fig:xgbtree}

% ================================ 5. RESULTS =================================
\section{Test Results}

\todo{All four models were evaluated once on the same untouched test set. Present the
metric table, then interpret the confusion matrices, ROC curves, and precision--recall
curves. Explain what each plot shows that accuracy alone would hide.}

\begin{table}[H]
\centering
\caption{Performance on the held-out test set.}
\label{tab:performance}
\small
\begin{tabular}{lrrrrrr}
\toprule
Model & Accuracy & Precision & Recall & F1 & ROC AUC & PR AUC \\
\midrule
Baseline (\todo{majority class}) & -- & -- & -- & -- & -- & -- \\
Decision tree & -- & -- & -- & -- & -- & -- \\
Random forest & -- & -- & -- & -- & -- & -- \\
AdaBoost & -- & -- & -- & -- & -- & -- \\
XGBoost & -- & -- & -- & -- & -- & -- \\
\bottomrule
\end{tabular}
\end{table}

\figslot{14_confusion_matrices}{0.96}{Confusion matrices for the four tuned models on the
test set.}{fig:confusion}
\figslot{15_roc_curves}{0.72}{ROC curves for the four tuned models.}{fig:roc}
\figslot{16_pr_curves}{0.72}{Precision--recall curves for the four tuned models.}{fig:pr}
\figslot{17_model_comparison}{0.82}{Test performance across models.}{fig:modelcomparison}

\subsection{Which method works best, and why}
\todo{Required by the assignment: compare the four models, say which works best, and give
your analysis of the reason (bias--variance, ensembling, boosting, regularization, data
size, noise, and so on). Say how large the differences are relative to uncertainty.}

% ======================= 6. LIMITATIONS & CONCLUSION =========================
\section{Assumptions, Limitations, and Conclusion}

\todo{List your assumptions, the dataset's limitations, anything the models cannot capture,
and a short conclusion that answers the objective stated in Section 1.}

% ================================ APPENDIX ===================================
\appendix
\section{Reproducibility and Requirement Check}

The notebook reads the raw data from the public URL below, so no local path is required.
The fixed random state is \todo{seed}. Scaling, feature selection, and tuning use training
data only. The notebook exports its experiment logs to the \texttt{results} directory.

\begin{center}
\small\url{\dataseturl}
\end{center}

\begin{table}[H]
\centering
\caption{Location of the required project components.}
\small
\begin{tabularx}{0.96\textwidth}{lX}
\toprule
Requirement & Evidence in the submitted work \\
\midrule
Public data source & \todo{URL and loading method} \\
Data consistency & \todo{missing, duplicate, and inconsistent-value checks} \\
Attribute summary and normality & \todo{summary table, plots, interpretation} \\
Categorical conversion & \todo{encoding used} \\
Standardization and normalization & \todo{both shown; which was used} \\
Correlation analysis & \todo{numerical matrix and plots} \\
Important-attribute selection & \todo{method and result} \\
Decision tree, random forest, AdaBoost, XGBoost & \todo{four tuned models} \\
Hyperparameter tuning & \todo{GridSearchCV grids and sizes} \\
Tree visualization & \todo{figures showing each model's trees} \\
Confusion matrix, precision, recall, F-score & \todo{table and figure} \\
ROC and precision--recall curves & \todo{figures and interpretation} \\
Model comparison & \todo{table and discussion} \\
Assumptions stated & \todo{section reference} \\
\bottomrule
\end{tabularx}
\end{table}

\bibliographystyle{plainnat}
\bibliography{references}

\end{document}
